\documentclass{rrparticle}
\usepackage{graphicx}

%Obviously, the commands below are not really needed when typesetting your contribution!!!
\newcommand{\miktex}{\hbox{Mik\kern-.15em\TeX}}
\newcommand{\texlive}{\hbox{\TeX Live}}

\title{Seismic moment tensor solutions for geophysical mapping of the Philippine Sea basin by GMT} 
\author[1]{Polina Lemenkova}

\affil[1]{Department Atomistilor 407, RO-077125, POB-MG6, Mgurele-Bucharest, Romania, EU.

Email: {\em polina.lemenkova@ulb.be}}


\keywords{Mapping, cartography, GMT, earthquake, geophysics, Philippine Sea, geology}
\pacs{01.30.-y, 01.30.Ww, 01.30.Xx}

\hyphenation{rrp-ar-ti-cle}

\begin{document}
\maketitle
\begin{abstract}
The study presents mapping seismic moment tensor solutions for analysis of geophysical settings on the margins of the Philippine Sea Plate (PSP). Active volcanism and seismic situation along the margins of the Philippine Sea Basin (PSB) shows earthquake events demonstrating high seismicity of the region caused by tectonic plate subduction. The aim was to perform cartographic mapping of the earthquake events using Generic Mapping Tools (GMT) modules, to analyze seismic situation through visualization of the geophysical datasets. The data was processed by GMT using multi-source data: GEBCO, CMT and ISC seismic data. Technically, a set of GMT modules was used ('psmeca', 'psvelo',  'grdtrack', 'grdimage') for data modelling and mapping. An application of GMT functionality is presented by selected code snippets with selected code snippets for 'psmeca' and 'psvelo', 'img2grd', 'img2grd' modules. Current paper presents a report on the cartographic techniques applied for geophysical mapping. Centroid moment tensor solutions are visualized for shallow depth earthquakes of Mw <10 along the PSB margins from 1976 to 2010. Earthquakes along the tectonic plate subduction zone are shown by GMT. Data from Global Centroid-Moment-Tensor (CMT) Project. Presented and explained GMT code snippets contribute to the development of methods in geological mapping using GMT scripting toolset, for visualizing focal mechanisms. 
\end{abstract}

\section{Introduction}

Geomorphological evolution of the oceanic trenches is largely controlled by a variety of factors including geological settings, tectonic plate movements and geomorphological profiling (Fujioka et al. 2002; Gong et al. 2017; Lemenkova, 2018; 2019a; 2019b), geophysical processes (Ogawa et al. 1997; Hall et al. 1995a; Seekings \& Teng, 1977), submarine volcanism and slope steepness (Ozawa et al. 2004; Lemenkova, 2019c, 2020f) inducing intensity of sedimentation and marine biological factors contributing to the sediment supply. 

Earlier studies on oceanic trenches highlighted relationship between the geomorphic patterns of the deep-sea channels, submarine fans and their topography related to the mid-ocean ridge tectonics, volcanism and dynamics of the Philippine back arc (Menard, 1955; Chang et al. 2007; Fujioka et al. 1999, Lemenkova, 2019d; Hall et al. 1995b), as well as earthquakes and gravitation (Lin \& Lo, 2013). Double seismic zone beneath the Mariana Island arc is explained (Samowitz \& Forsyth, 1981) by the following conceptual scheme of the processes of the global plate tectonics. Subduction of the cooled plate into the mantle causes formation of the deep ocean trenches where, as a consequence, earthquakes and tsunamis originate. 

Active volcanism and seismic situation around Philippine Sea margins (Fig. 5) shows earthquake events demonstrating high seismicity of the region caused by tectonic plates subduction. Subducting lithospheric plates act as giant radiators of the heat cooling, thickening, and progressively subsiding from ridge to trench. As a result, spreading seafloor, affected by moving plates, creates an axial rift, corrugated hills, and ridges, formed by the nearby faults. Plate subduction along the Izu-Bonin-Mariana (IBM) trench caused trench roll-back, arc rupture and back-arc rifting of two interconnected back-arc basins: Parece Vela and Shikoku Basins. The KPR, a remnant arc of the active IBM system results from the spreading of the Shikoku Basin (Fig. 1). It plays a key role in the subduction process of the west Philippine Sea Plate: the subduction zone here is characterized by the Kyushu-Palau Ridge (KPR) subducting beneath the Kyushu Arc, Fig. 1 (Cao et al. 2014). 

A set of geodynamic processes including subduction-interface rheology, phase-transition buoyancy, slab stagnation, rollback of mantle and ridge-push effects, cause significant trench motion of the IBM detected as advance by global plate-motion observations (Cížková, H., Bina, 2015). 

\section{Materials and Methods}
\subsection{Seismic mapping}

Focal mechanisms (Fig. 2) were drawn using 'psmeca' module of Generic Mapping Tools (GMT) that reads dataset values from ASCII file and generates a PostScript code plotting focal mechanisms. Focal mechanisms depict a commonly accepted theory of the focal depths of the earthquake, showing a depth from the surface to the earthquake's origin (a.k.a. hypocenter). Given dataset contains significant amount of shallow earthquakes of the PSB with focal depths of several tens km (mostly $<50$), intermediate earthquakes with focal depths from 70 to 200 km, and also a few earthquakes of PSB with deep focus reaching depths $> 500$ km, Table tab:1. The foci of the most PSB earthquakes are concentrated in the crust and upper mantle, that is, originate in shallow parts of the Earth's interior.

The original dataset of the focal mechanism solutions are from the was taken from the Global CMT Catalog, formerly known as the Harvard CMT catalog (Ekström et al. 2012). Technical parameters are set up as following: 
\begin{enumerate}\itemsep 0em 
  \item Moment magnitude (Mw) $<10$; 
  \item Surface wave magnitude (Ms): $<10$);
  \item Body wave magnitude (Mb): $<10$;  
  \item Data time span: 1976-2010
 \end{enumerate}
	While magnitude is a widely understood concept, describing the energy of the earthquake release on a logarithmic scale, some technical details are presented in Table 1 and Table 2 used for mapping velocity (Fig. 4), assigning and interpreting magnitudes. Introduced in the 1930's, the Richter Scale is the best known scale for measuring the magnitude of earthquakes (Müller, 2001). Surface wave magnitude (Ms), creating the strongest disturbance within the upper layers of the Earth, is computed by the following algebraic formula:
	$MS=log10(A/T)+1.66log10(D)+3.30$,
where T is the measured wave period and D is the distance in radians. 
	
	On the contrary, earthquakes occurring deep in the Earth do not generate large surface waves. Therefore, a body wave magnitude (Mb) is scaled based on the seismic waves penetrating through the Earth's interior (body) (Müller, 2001). Body wave magnitude (mb) is measured based on the maximum amplitude A by formula:
	$Mb=log10(A/T)+Q(D,h)$,
where T is the measured wave period and Q is an empirical function of focal depth h and epicentral distance D.
	
	Moment magnitude (MW) relies on an underlying robust physical and mathematical development through converting seismic moment using a formula calibrated to agree with MS over much of its range. Others parameters are taken as default ones. Data are collected in both GMT 'psvelomeca' and GMT 'psmeca' input formats and tested. Final solution was for GMT 'psmeca' (the reason is for the GMT version 5.4.5 compatibility). The focal mechanisms are derived from a solution of the earthquakes' moment tensor, which is estimated by an analysis of the observed seismic waveforms. 
	
	Technically, seismic moment tensor (Harvard CMT, with zero trace) was plotted (Fig. 2) by 'psmeca' module using the following GMT code snippet: $gmt psmeca -R CMT.txt -J -Sd0.5/8/u -Gred -L0.1p -Fa/5p/it -Fepurple -Fgmagenta -Ft -W0.1p -Fz -Ewhite -O -K >> \$ps$.
	
	The scale (0.5) was adjusted to the scaling of the 'beach ball' radius, which is proportional to the magnitude. Here the '-Sd0.5/8/u' implies the scale, label size and annotation placement below the 'beach ball'. Filling of the extensive quadrants was defined by -E parameter; -L parameter defines drawing the 'beach ball' outline of 0.1 pt. Shaded compressional quadrants of the focal mechanism 'beach ball' are colored by -Gred function. The -Fa/5p/it option was used to plot size, P\_axis\_symbol and T\_axis\_symbol, to compute and plot P and T axes with symbols (here: selected inverse triangle (i) and triangle (t), respectively). 

\section{Specific Commands}

The present style modifies a couple of common \LaTeX ~commands and implements a few new ones. The modified commands are related mostly to the layout on the first page of the document (title, authors, affiliations and abstract).

The \texttt{\small \textbackslash title} command accepts the new line character ''\textbackslash\textbackslash`` in its argument but, when inserted into the document, the title is written uppercase. The same transformation is forced upon first and second level section titles (\texttt{\small \textbackslash section} and \texttt{\small \textbackslash subsection} commands). There is however a special command to be used whenever acronyms or words with special lowercase letters must appear in the text, namely  \texttt{\small \textbackslash NoCaseChange\{\}} (defined in the \textbf{textcase} package). Its effect is that the argument appears in the text unmodified by letter casing commands (also implemented by \textbf{textcase} package). For instance
\begin{center}
\texttt{\small
\textbackslash title\{Lagrangean \textbackslash NoCaseChange\{sp(3)\} BRST Formalism for Massive Vectorial Bosonic Fields\}
}
\end{center}
will appear on the first page as
\begin{center}
LAGRANGEAN sp(3) BRST FORMALISM FOR MASSIVE VECTORIAL BOSONIC FIELDS.
\end{center}

The \textbackslash author command is slightly different than the one implemented by the standard \textit{article} class. Its new syntax is
\begin{center}
\texttt{\small
\textbackslash author[\textit{affilSign}]\{\textit{authorName}\}
}
\end{center}
and must be issued for each of the authors. It adds the \textit{authorName} at the end of the list of authors for your paper (beware that the order of the authors in the list is therefore given by the order of their corresponding \textbackslash author commands!). It also assigns the affiliation key (or list of comma separated keys) \textit{affilSign} to the current author. The keys are recommended to be either number or lower case letters and they should be set incrementally starting with '1', respectively 'a'. In the case that all the authors have the same affiliation this key may be absent. As one may notice analyzing the source of this document, we prefer the contact information such as emails to appear in the affiliation block corresponding to the author, alphabetically indexed (whenever the affiliation is shared with other author(s)) and in italics. \textbf{Beware} that the responsibility of using the appropriate keys for author(s) and affiliation(s) belongs to the authors of the scientific material! Obvious inconsistencies will be scrutinized by the editorial personnel when encountered but the RRP editorial office will not check the validity of the provided data nor it can be held responsible if invalid affiliations are published! As one may also notice the first argument of the \textbackslash author command can contain links to various notes regarding the status of a specific author and this information is introduced using the command 
\begin{center}
\texttt{\small
\textbackslash authnote\{\textit{noteText}\}
}
\end{center}
(see the \LaTeX ~source of this document for an example -- Jane Doe). No comma must appear after the last \textit{affilSign} key and a following \textbackslash authnote command.


In close relation to the author list, one must introduce the affiliation list using the command
\begin{center}
\texttt{\small
\textbackslash affil[\textit{affilSign}]\{\textit{affiliationText}\},
}
\end{center}
where \textit{affilSign} is the specific symbol key used for (and only for) the currently entered affiliation and which must also appear next to at least one of the authors in the author list. As for the author list the affiliation list must be defined in the exact order they need to appear in the text (keeping in mind that the keys must be in low-to-high order). The affiliation key may again be missing if and only if all the authors are affiliated to the same department, institute/company and so on. The \textit{affiliationText} may contain the end-of-line command that should be used at specific points in the text where a new line is required for aesthetic reasons(symmetry of centered lines) or clarity. We recommend however that one should not abuse of this feature in order to minimize the vertical space occupied by such informations on the first page.

Among the new commands defined by the RRP style there are 
\begin{center}
\texttt{\small
\textbackslash keywords\{\textit{comma separated key words list}\}
}\end{center}
which \textbf{must} appear in the header of the document (before \textbackslash begin\{document\}) or at most before \textbackslash begin\{abstract\} is issued to have any effect at all. Moreover the list of key words should not end with a period '.' as it is automatically added by the class code.
\begin{center}
\texttt{\small
\textbackslash pacs\{\textit{PACS number list}\}
}\end{center}
allows one to specify PAC identifiers as a comma separated list.

There is a special environment for specifying the acknowledgments
\begin{center}
\texttt{\small
\textbackslash begin\{acknowledgement\}...\textbackslash end\{acknowledgement\}
}\end{center}
The acknowledgments are usually typeset in a font face of size 9pt so please do not use font size changing commands (they will be removed in editorial processing). Special formatting such as bold, italic or slanted are accepted and recommended as the way to emphasize any fragments of the text that you consider necessary.

\section{Examples for various Environments}

This section includes examples of different environments containing media and data material (the copyright of which is already owned by RRP) much needed in any good scientific publication. For instance below (in figure \ref{pic1}) we present a \textit{figure} environment as it must appear in every contribution submitted to our journal (this picture is in PNG format so compiling must be done using \textit{pdflatex} command rather than the usual \textit{latex} command or one should specify the \textbf{pdftex} driver when loading the \textbf{graphicx} package).
\begin{figure}[h!tb]
\centering
\includegraphics[width=0.8\textwidth]{xNLS_DAMI.png}
\caption{This is a sample picture taken from RRP volume \textbf{50}(1-2) from page 129 (2005).It illustrates (hashed areas) the instability domains for a series of nonlinear Schr\"odinger equations determined using the deterministic approach to modulational instability.}
\label{pic1}
\end{figure}

Another environment with special formatting in our style is the \textit{table} environment, a sample of which is table \ref{table1}.
\begin{table}[h!t]%
\caption{This table is taken from RRP volume \textbf{50}(1-2) from page 43 (2005). It gives the ``
\textit{number of bound states dependence on the radius of space curvature for $\alpha = 0.005$, $U_0 = 1$}''.}
\centering
\begin{tabular}{|c|c|}
\hline
Value $\rho$ & Value $\varepsilon$ \cr
\hline
$\rho = 50$ & -- \cr
$\rho = 100$ & -- \cr
$\rho = 250$ & $\varepsilon_1 = 0.0289$ \cr
$\rho = 400$ & $\varepsilon_1 = 0.3772$ \cr
$\rho = 1000$ & $\varepsilon_1 = 0.4142$, $\varepsilon_2 = 0.8495$ \cr
\hline
\end{tabular}
\label{table1}
\end{table}

As RRP is a physics journal, an important part of the scientific language used in the publication is constituted by equations. One of the main reason for which the developer decided to include the AMS style within the RRP style is the appropriate formatting of the mathematical environments. For instance the \textit{split} environment is recommended as a very good replacement of the \textit{eqnarray} environment. You can see the difference between \textit{split} and \textit{eqnarray} in equation \eqref{eq1} and \eqref{eq2} respectively
\begin{equation}\label{eq1}\begin{split}
S_{0}&[A^{\mu },\dot{A}^\mu,\phi ]=\int d^{4}x\mathcal{L}=\int
d^{4}x\ (-\frac{1}{4}F_{\mu \nu }F^{\mu \nu }- \\
&-k~\partial _{\lambda }F^{\alpha \lambda }\partial _{\rho }F_{\alpha }^{\rho
}+\frac{1}{2}(\partial _{\mu }\phi -mA_{\mu })(\partial ^{\mu }\phi -mA^{\mu
}))
\end{split}\end{equation}
\begin{eqnarray}
& &S_{0}[A^{\mu },\dot{A}^\mu,\phi ]=\int d^{4}x\mathcal{L}=\int
d^{4}x\ (-\frac{1}{4}F_{\mu \nu }F^{\mu \nu }- \label{eq2} \\
& &-k~\partial _{\lambda }F^{\alpha \lambda }\partial _{\rho }F_{\alpha }^{\rho
}+\frac{1}{2}(\partial _{\mu }\phi -mA_{\mu })(\partial ^{\mu }\phi -mA^{\mu
})) \nonumber
\end{eqnarray}
The AMS style also allows one to easily make use of subequations
\begin{subequations}
\label{ein}
\begin{align}
\frac{\ddot a_2}{a_2} +\frac{\ddot a_3}{a_3} +\frac{\dot
a_2}{a_2}\frac{\dot
a_3}{a_3} - \frac{n^2}{a_3^2} &= \kappa T_{1}^{1}\,, \label{11}\\
\frac{\ddot a_3}{a_3} +\frac{\ddot a_1}{a_1} +\frac{\dot
a_3}{a_3}\frac{\dot
a_1}{a_1} - \frac{m^2}{a_3^2} &= \kappa T_{2}^{2}\,, \label{22} \\
\frac{\ddot a_1}{a_1} +\frac{\ddot a_2}{a_2} +\frac{\dot
a_1}{a_1}\frac{\dot
a_2}{a_2} + \frac{m n}{a_3^2} &= \kappa T_{3}^{3}\,, \label{33} \\
\frac{\dot a_1}{a_1}\frac{\dot a_2}{a_2} +\frac{\dot a_2}{a_2}
\frac{\dot a_3}{a_3} + \frac{\dot a_3}{a_3}\frac{\dot a_1}{a_1} -
\frac{m^2 - m n + n^2}{a_3^2} &=
\kappa T_{0}^{0}\,, \label{00}\\
m \frac{\dot a_1}{a_1} - n \frac{\dot a_2}{a_2} - (m - n) \frac{\dot
a_3}{a_3} &= \kappa T_{3}^{0}\,. \label{03}
\end{align}
\end{subequations}
though one should consult the AMS user guide \cite{amsug} when special features like splitting subequations between subsequent pages are needed (see \cite{amsug} ver. 2.0, revised 2002, p. 9-10). Please notice that the \textit{subequations} environment presents the advantage of referencing both the whole group of equations \eqref{ein} or each of the sub-equations individually, \textit{e. g.}  \eqref{33}. Of course, there are other mathematical AMS environments that the author may use in their papers and to learn about them one should consult the afore-mentioned user guide.

\section{Instead of Conclusions}

The current style provides a few elements of both authenticity for us (the headers and footers) as well as a way for you, the author, to certify (to third parties) the submission of a manuscript to our journal. Most of the formatting is done automatically in the background so you \textbf{must not} interfere with elements such as page layout parameters, spacing parameters used in various environments or font sizes. We admit however that in certain conditions (especially when lots of floating environments (figures and tables) are used) the author may need to issue special commands to force the output of floats on specific pages (such commands are \textbackslash clearpage, \textbackslash newpage, etc.). Of course these commands are allowed in your \LaTeX ~code. We also recommend that you do a spelling before submitting your article and use the \textbackslash hyphenation command (or hyphenation marks) whenever necessary.

These recommendations should be taken into account together with the instructions for authors available for consulting on the Romanian Reports in Physics web site when submitting articles written in \LaTeX ~to our journal. Please, be advised that failure to comply may lead to undefinite delays in the publication of your submitted article even if the referees gave you a positive recommandation!

The current style (Dec. 2010) is however not a very well tested version, therefore comments and contributions to the development of the code as well as help in debugging specific issues is and will always be highly regarded by our editorial team!  

\begin{acknowledgement}
The author(s) would like to acknowledge the help received from Mrs. Margareta Oancea and the members of the technical department at the Romanian Academy Publishing House, in principal, in specifying a \LaTeX format that matches as close as possible the MS Word template that RRP uses officially, but also for carefully and patiently correcting the various versions produced during development. Last but not least, I owe gratitude to those who made possible that Donald Knuth's \cite{dknuthhp} complex but rigorous \TeX \cite{knuth,teximpacient} evolve into the friendlier \LaTeXe (of course names like Leslie Lamport \cite{leslie} or the people in the \LaTeX3 project led by Frank Mittelbach \cite{mittelbach,ltexclass,ltexshort,ltexclass2} come to mind). A big ``Thanks!'' goes to all the web sites, forums and message lists maintainers out there on the internet for the documentation and examples they shared with the world in the obvious effort of teaching people about the ways and wonders of beautiful typesetting (a small part of these resources are \cite{texrefex,eijkhout,texcookbook,pagelayout,tug,stacko,web1,web2,web3,web4,web5,web6,web7,web8}). 
\end{acknowledgement}

\begin{thebibliography}{99}
\bibitem{knuth} D. E. Knuth, D. R. Bibby, ``\textit{The \TeX book}'' 20th edn. (AMS \& Addison-Wesley Publ. Co., 1991)
\bibitem{dknuthhp} D. E. Knuth homepage: \texttt{\small http://www-cs-faculty.stanford.edu/\~{}knuth/}
\bibitem{eijkhout} V. Eijkhout, ``\textit{\TeX ~by Topic, A \TeX nician's Reference}'' (electronic ver. 1.0, 2001)
\bibitem{texcookbook} ``\textit{\TeX Cookbook}'' (MathPro Press, 1989; \texttt{\small http://www.mathpro.com/math/MathPro.html})
\bibitem{texrefex} D. Bausum, ``\textit{\TeX: Reference and Examples}'' (\texttt{\small http://www.tug.org/utilities/ plain/cseq.html})
\bibitem{teximpacient} P. W. Abrahams, K. A. Hargreaves, K. Berry, ``\textit{\TeX for the Impatient}'' (Addison-Wesley Publ. Co., 2003)
\bibitem{tug}\TeX  \textit{Users Group}: \texttt{\small http://www.tug.org/tutorials/}
\bibitem{leslie} L. Lamport, ``\textit{\LaTeX. A Document Preparation System}'' 2nd edn. (Addison-Wesley Publ. Co., 1994)
\bibitem{mittelbach} F. Mittelbach, M. Goossens (J. Braams, D. Carlisle, C. Rowley), ``\textit{The \LaTeX Companion}'' 2nd edn. (Addison-Wesley, Boston, 2004)
\bibitem{ltexclass} The \LaTeX3 Project, ``\textit{\LaTeXe for Class and Package Writers}'' (1999)
\bibitem{ltexshort} T. Oetiker, H. Partl, I. Hyna, E. Schlegl, ``\textit{The Not So Short Introduction to \LaTeXe}'' (ver. 4.31, June 24, 2010)
\bibitem{ltexclass2} The \LaTeX3 Project, ``\textit{\LaTeXe for Authors}'' (2001)
\bibitem{pagelayout} P. van Oostrum, ``\textit{Page Layout in \LaTeX}'' (\texttt{\small http://www.cs.ruu.nl/people/piet})
\bibitem{web1}A. J. Hildebrand, \texttt{\small http://www.math.uiuc.edu/\~{}hildebr/tex/}
\bibitem{web2} \texttt{\small http://zoonek.free.fr/LaTeX/}
\bibitem{web3} \texttt{\small http://amath.colorado.edu/documentation/LaTeX}
\bibitem{web4} \texttt{\small http://www-h.eng.cam.ac.uk/help/tpl/textprocessing/}
\bibitem{web5} \texttt{\small http://latex.computersci.org/}
\bibitem{web6} S. Kottwitz, \texttt{\small http://texblog.net/}
\bibitem{web7} \texttt{\small http://www.personal.ceu.hu/tex/}
\bibitem{web8} The UK \TeX Archive: \texttt{\small http://www.tex.ac.uk/}
\bibitem{stacko} Stack Overflow: \texttt{\small http://stackoverflow.com/tags}
\bibitem{amsug} User's Guide for the \textit{amsmath} Package: \texttt{\small ftp://ftp.ams.org/pub/tex/doc /amsmath `/amsldoc.pdf}
\end{thebibliography}

\end{document}